We formalize trajectory selection as a \emph{hierarchical} decision problem: given a set of diverse candidates from a multi-modal generator, the objective is to select the single most acceptable ego trajectory under a prioritized rulebook.  We move beyond treating selection as a minor post-processing step, instead encoding the non-negotiable priority structure of traffic rules through tiered, tolerance-based lexicographic ordering.  The resulting formulation provides a clean interface between (i)~any multi-modal generator
and (ii)~an explicit, auditable rulebook evaluator.

The development proceeds in a deliberate sequence.  We first fix notation and define the prediction setting (Sections~\ref{subsec:notation}--\ref{subsec:prediction_setting}), then introduce the geometric metrics that characterize candidate-set quality and show why they are insufficient for acceptability (Sections~\ref{subsec:geometric_metrics}--\ref{subsec:error_distribution}). With this motivation in place, we formalize the rule taxonomy, tier-score aggregation, and the lexicographic selection criterion (Sections~\ref{subsec:rule_taxonomy}--\ref{subsec:lexicographic_selection}), establish its key properties (Section~\ref{subsec:properties}), and verify three correctness properties (Section~\ref{subsec:formal_guarantees}).  We close with a comparison to alternative strategies (Section~\ref{subsec:strategy_comparison}), the dual-protocol evaluation context (Section~\ref{subsec:fullcatalog_context}), and the formal problem
statement (Section~\ref{subsec:problem_statement}).

%-------------------------------------------------------------------------------
\subsection{Notation and Conventions}
\label{subsec:notation}
%-------------------------------------------------------------------------------

Table~\ref{tab:notation} summarizes the notation used throughout the paper. We follow a scenario-centric viewpoint: each scenario~$s$ contains ego history, nearby agents, and map context, and the predictor produces a candidate set~$\mathcal{T}$ along with confidence scores.  Rules are organized into four tiers, and violations are aggregated per tier.  All symbols introduced in this table are used consistently in the sections that follow; we reference them without redefining them.

\begin{table*}
\centering
\small
\caption{Key Notation and Symbol Definitions}
\label{tab:notation}
\begin{tabular}{ll}
\toprule
\multicolumn{2}{l}{\textit{Scenarios and Trajectories}} \\
$s$ & Driving scenario sampled from dataset distribution $\mathcal{D}$ \\
$\mathbf{x}_{\text{ego}}$ & Ego history (past states over $T_h$ timesteps) \\
$\mathbf{X}_{\text{agents}}$ & Neighbor-agent histories \\\
$\mathcal{M}$ & Map context (lanes, signals, crosswalks, boundaries, etc.) \\
$\tau$ & Ego future trajectory over $T$ timesteps \\
$\tau_t$ & Ego state at timestep $t$ (typically position and heading) \\
$\mathcal{T}$ & Set of $K$ candidate trajectories \\
$\tau^{\mathrm{gt}}$ & Ground truth future trajectory \\
$\Delta t$ & Sampling interval (0.1\,s at 10\,Hz) \\
\midrule
\multicolumn{2}{l}{\textit{Confidence and Modes}} \\
$p^{(k)}$ & Confidence score for candidate $k$ \\
$K$ & Number of candidate modes (typically 6) \\
$T_h$ & History length (11 timesteps, 1.1\,s) \\
$T$, $T_f$ & Prediction horizon (50 timesteps, 5.0\,s); $T_f$ used in architecture descriptions \\
$N_{\text{scen}}$ & Number of scenarios in a scenario-level aggregation (12{,}800 in WOMD \texttt{validation\_interactive}) \\
$N_{\text{eval}}$ & Number of evaluation samples (one per scenario in the current pipeline; $N_{\text{eval}} = N_{\text{scen}}$) \\
\midrule
\multicolumn{2}{l}{\textit{Rules and Compliance}} \\
$\mathcal{R}$ & Rule catalog; $R = |\mathcal{R}| = 28$ \\
$r$ & Individual traffic rule \\
$\ell(r)$ & Tier assignment for rule $r$ in $\{0,1,2,3\}$ \\
$a_r(s)$ & Oracle applicability indicator (auditing): 1 if rule $r$ applies in scenario $s$ \\
$\hat{a}_r(s)$ & Predicted applicability used by RECTOR for selection (thresholded from model output) \\
$\tilde{a}_r(s)$ & Continuous applicability score before thresholding; $\tilde{a}_r = \sigma(\text{logit}_r)$ \\
$V_r(\tau; s)$ & Raw violation severity of $\tau$ under rule $r$ in scenario $s$ \\
$\bar{V}_r(\tau; s)$ & Normalized violation severity (mapped and clamped to $[0,1]$) \\\
$\alpha_r$ & Per-rule normalization constant; $\bar{V}_r = \mathrm{clamp}(V_r / \alpha_r,\,0,\,1)$ \\
\midrule
\multicolumn{2}{l}{\textit{Tier Scores and Selection}} \\
$\tilde{w}_r$ & Intra-tier weight for rule $r$; $\sum_{r:\,\ell(r)=\ell}\tilde{w}_r=1$ per tier (default: uniform) \\
$S_\ell(\tau; s)$ & Aggregated tier-$\ell$ score \\
$\mathbf{S}(\tau; s)$ & Tier score vector $[S_0,S_1,S_2,S_3]$ \\
$\epsilon_\ell$ & Tier tolerance used for robust lexicographic comparison \\
$\prec_{\text{lex},\epsilon}$ & Shorthand for procedural lexicographic selection (Definition~\ref{def:lex_selection}); not a partial order \\
$L$ & Number of tiers ($L=4$) \\
\midrule
\multicolumn{2}{l}{\textit{Metrics}} \\
$\text{minADE}$ & Minimum Average Displacement Error over modes \\
$\text{minFDE}$ & Minimum Final Displacement Error over modes \\
$\text{MissRate}$ & Fraction of scenarios with $\text{minFDE} > \delta$ \\
\bottomrule
\end{tabular}
\end{table*}

%-------------------------------------------------------------------------------
\subsection{Trajectory Prediction Setting}
\label{subsec:prediction_setting}
%-------------------------------------------------------------------------------

With the notation in place, we can define the prediction setting precisely. Let $s \sim \mathcal{D}$ denote a driving scenario drawn from a dataset distribution.  A multi-modal predictor with parameters~$\theta$ maps each scenario to (i)~a finite set of candidate ego trajectories and (ii)~a
confidence score for each candidate:
\begin{align}
(\mathcal{T}, \mathbf{p}) = f_\theta(s), \quad
\mathcal{T} = \{\tau^{(1)},\ldots,\tau^{(K)}\}, \quad
\mathbf{p} = (p^{(1)}, \ldots, p^{(K)}),\nonumber
\end{align}
where each candidate $\tau^{(k)}$ is a length-$T$ sequence of ego states at sampling interval~$\Delta t$.  The goal is not only to generate diverse candidates, but also to select a single trajectory~$\hat{\tau}$ that is preferable under a prioritized rulebook.  The remainder of this section
formalizes what ``preferable means.

%-------------------------------------------------------------------------------
\subsection{Standard Geometric Metrics}
\label{subsec:geometric_metrics}
%-------------------------------------------------------------------------------

Before introducing the rule-based selection objective, we define the standard geometric metrics used in multi-modal trajectory prediction, because understanding their limitations motivates the formulation that follows. These metrics characterize the \emph{coverage} quality of the candidate set~$\mathcal{T}$, but---crucially---they do not determine whether the \emph{selected} mode is acceptable under traffic rules.

\begin{definition}[Average Displacement Error] The Average Displacement Error (ADE) is the mean Euclidean distance over the prediction horizon:
\begin{equation}
\text{ADE}(\tau, \tau^{\mathrm{gt}}) = \frac{1}{T} \sum_{t=1}^T \| \tau_t - \tau^{\mathrm{gt}}_t \|_2.
\end{equation}
\end{definition}

\begin{definition}[Final Displacement Error] The Final Displacement Error (FDE) is the Euclidean distance at the final timestep:
\begin{equation}
\text{FDE}(\tau, \tau^{\mathrm{gt}}) = \| \tau_T - \tau^{\mathrm{gt}}_T \|_2.
\end{equation}
\end{definition}

\noindent
For multi-modal prediction, we report the minimum error across candidates, which measures the best achievable match to the ground truth
within~$\mathcal{T}$:
\begin{align}
\text{minADE}(s) &= \min_{k \in [K]} \text{ADE}(\tau^{(k)}, \tau^{\mathrm{gt}}), \\
\text{minFDE}(s) &= \min_{k \in [K]} \text{FDE}(\tau^{(k)}, \tau^{\mathrm{gt}}).
\end{align}

\begin{definition}[Miss Rate]
The miss rate is the fraction of scenarios where the best candidate exceeds a threshold~$\delta$:
\begin{equation}
\text{MissRate}_\delta = \mathbb{E}_{s \sim \mathcal{D}} \left[ \mathbb{1}\big[\text{minFDE}(s) > \delta\big] \right].
\end{equation}
\end{definition}

\noindent
We use $\delta = 2.0\,\mathrm{m}$ following the Waymo benchmark convention. MissRate does not enter the lexicographic selection objective; it is included here as a standard benchmark metric that characterizes candidate-set coverage.  The actual motivation for rule-aware selection comes from the heavy-tailed error analysis in Section~\ref{subsec:error_distribution}, which uses percentile statistics rather than a single miss-rate threshold. In addition to these best-of-$K$ metrics, our evaluation in later sections reports \emph{selected-trajectory} metrics (selADE/selFDE), which quantify the quality of decisions at the selection stage rather than the coverage of the candidate set.

%-------------------------------------------------------------------------------
\subsection{Error Distribution Analysis}
\label{subsec:error_distribution}
%-------------------------------------------------------------------------------

The metrics above reveal a critical property of the prediction setting that directly motivates rule-aware selection: the distribution of prediction errors is \emph{heavy-tailed} in interactive driving.  This matters because selection mistakes in the tail can lead to safety- and law-critical failures, even when median performance is strong.  Using scenario-level aggregation over $N_{\text{scen}}{=}12{,}800$ interactive WOMD scenarios (as distinct from the canonical sample-level count $N_{\text{eval}}{=}25{,}448$ defined in Table~\ref{tab:notation} and Section~\ref{Sec:Experiments}), the per-scenario minADE confirms this: the median is 1.05\,m, but the 90th percentile exceeds 8\,m, the 95th reaches 10.47\,m, and the 99th reaches 16.55\,m---roughly $16\times$ the median (Fig.~\ref{fig:percentile_analysis}).

In such cases, confidence-only selection is especially brittle, because mode probabilities can be poorly calibrated under rare interactions, while rule constraints (e.g., collision margins or red-light compliance) remain decisive for acceptability.  This motivates a selection mechanism grounded in traffic rules rather than likelihood alone.
The formal machinery for such a mechanism is developed next.

%-------------------------------------------------------------------------------
\subsection{Traffic Rule Taxonomy}
\label{subsec:rule_taxonomy}
%-------------------------------------------------------------------------------

To define rule-aware selection, we must first specify how traffic rules are represented.  Let $\mathcal{R}$ denote a catalog of $R{=}28$ traffic rules. Each rule $r \in \mathcal{R}$ is assigned to one of four tiers via $\ell(r)\in\{0,1,2,3\}$, reflecting the normative priority structure of
traffic constraints:
\begin{equation}
\text{Safety} \succ \text{Legal} \succ \text{Road} \succ \text{Comfort},
\end{equation}
where higher tiers must dominate lower tiers in the selection decision (see Section~\ref{subsec:observations} for the normative justification and the Legal-tier disclaimer).

Rules are not universally applicable---a stop-sign rule is irrelevant in a scene without a stop sign---and this context dependence is one of the central technical challenges addressed in this work.  Two distinct applicability indicators are used throughout the paper, and the distinction between them is
\emph{operationally critical}:
\begin{itemize}[leftmargin=1.5em]
\item \textbf{Oracle applicability} $a_r(s)\in\{0,1\}$: a ground-truth   label derived from map and scenario metadata.  It is used \emph{exclusively} for post-hoc auditing (Protocol~B in Section~\ref{subsec:fullcatalog_context}) and \emph{never} enters the selection objective.
\item \textbf{Predicted applicability} $\hat{a}_r(s)\in\{0,1\}$: the
  thresholded output of a learned neural head (Section~\ref{Sec:RECTOR}) that
  predicts rule relevance from the scene embedding.  This is the indicator
  used in the tier score aggregation (Definition~\ref{def:tier_score}) and in
  the selection objective (Problem~\ref{prob:selection}).
\end{itemize}
\noindent
Oracle applicability is defined as:
\begin{equation}
a_r(s)=
\begin{cases}
1 & \text{if rule $r$ applies in scenario $s$},\\
0 & \text{otherwise.}
\end{cases}
\end{equation}
The predicted counterpart $\hat{a}_r(s)$ is obtained by thresholding a continuous score $\tilde{a}_r(s) = \sigma(\text{logit}_r)$ produced by the applicability head.  The gap between $a_r(s)$ and $\hat{a}_r(s)$ is the applicability prediction error, and it is one of the key sources of approximation in the deployed system.  This motivates the dual-protocol evaluation design (Section~\ref{subsec:fullcatalog_context}).

Each applicable rule produces a non-negative violation severity $V_r(\tau;s)$ for a candidate trajectory~$\tau$ in scenario~$s$:
\begin{equation}
V_r(\tau;s)\ge 0.
\end{equation}
Severity can be binary or continuous (e.g., margin-based), and it can depend on kinematics, interaction context, and map semantics.  Because the 28~rules are heterogeneous---their raw severities carry different units (metres for collision margins, m/s for speed violations, dimensionless IoU for overlap
rules)---a \emph{normalization contract} is required before tier-level aggregation.  We define the normalized violation as 
\begin{equation}
\label{eq:normalization_contract}
\bar{V}_r(\tau;s) \;=\; \mathrm{clamp}\!\Big(\frac{V_r(\tau;s)}{\alpha_r},\;0,\;1\Big),
\end{equation}
where $\alpha_r > 0$ is a per-rule normalization constant that maps the physically meaningful range of~$V_r$ to the unit interval. This ensures $\bar{V}_r(\tau;s)\in[0,1]$ for every rule, and consequently each tier score $S_\ell(\tau;s)\in[0,1]$ (Definition~\ref{def:tier_score}). The specific values of~$\alpha_r$ and the clamping details for each rule are given in Section~\ref{Sec:Rule_Evaluation_And_Data_Pipeline}; only the
contract $\bar{V}_r\in[0,1]$ is needed for the selection formulation that follows.  This separation is deliberate: the present section defines the \emph{selection objective}, while Section~\ref{Sec:Rule_Evaluation_And_Data_Pipeline} defines the \emph{concrete evaluators} that instantiate it.

%-------------------------------------------------------------------------------
\subsection{Tier Score Aggregation}
\label{subsec:tier_aggregation}
%-------------------------------------------------------------------------------

Given the per-rule violations and applicability indicators defined above, the next step is to aggregate them into tier-level scores for comparison by the lexicographic selector.  For each tier $\ell \in \{0,1,2,3\}$, we aggregate violations across the rules assigned to that tier, using \emph{predicted} applicability $\hat{a}_r(s)$ for gating during selection (not the oracle indicator, which is reserved for auditing).

\begin{definition}[Tier Score]
\label{def:tier_score}
The tier-$\ell$ score of trajectory $\tau$ in scenario $s$ is
\begin{equation}
S_\ell(\tau;s)=\sum_{r \in \mathcal{R}:\,\ell(r)=\ell} \tilde{w}_r\,\hat{a}_r(s)\,\bar{V}_r(\tau;s),
\end{equation}
where $\tilde{w}_r \geq 0$ are per-rule weights satisfying $\sum_{r:\,\ell(r)=\ell}\tilde{w}_r = 1$ within each tier (defaulting to uniform: $\tilde{w}_r = 1/n_\ell$, where $n_\ell$ is the number of rules in tier~$\ell$).  This normalization ensures $S_\ell(\tau;s)\in[0,1]$ for all $\ell$, regardless of the number of rules per tier---a property required by Proposition~\ref{thm:scalarization}.
\end{definition}

\noindent
The four tier scores are collected into a tier score vector:
\begin{equation}
\mathbf{S}(\tau;s)=\big[S_0(\tau;s),\,S_1(\tau;s),\,S_2(\tau;s),\,S_3(\tau;s)\big]\in\mathbb{R}_{\ge 0}^{4}.
\end{equation}

This construction has four practical properties that make it suitable for hierarchical selection.  First, it is non-negative and bounded in $[0,1]$ per tier, which supports stable comparisons across candidates and guarantees scalarization correctness (Proposition~\ref{thm:scalarization}). Second, the intra-tier normalization ensures that a tier with 14~rules (Comfort) does not numerically dominate a tier with 2~rules (Road) at equal per-rule severity.  Third, cross-tier priorities are represented \emph{explicitly} in the tier structure rather than \emph{implicitly} in a weight vector.  Fourth, applicability gating prevents irrelevant rules from dominating the score when context is missing (e.g., stop-sign constraints in scenes without stop signs).

%-------------------------------------------------------------------------------
\subsection{Lexicographic Selection Criterion}
\label{subsec:lexicographic_selection}
%-------------------------------------------------------------------------------

With tier scores defined, we now specify how they are compared across candidates to enforce the priority structure.  The core mechanism is lexicographic ordering over the tier score vector.  Because numeric scores can be close among candidates, we use a tolerance-based comparator to avoid instability due to floating-point noise and near-ties.

\begin{definition}[Tolerance-Based Lexicographic Selection (Procedural)]
\label{def:lex_selection}
Given candidates $\mathcal{T}=\{\tau^{(1)},\ldots,\tau^{(K)}\}$, tier score
vectors $\mathbf{S}^{(k)}=\mathbf{S}(\tau^{(k)};s)$, and per-tier
tolerances $(\varepsilon_0,\ldots,\varepsilon_3)$ with $\varepsilon_\ell\ge 0$,
the selection proceeds by iterative pruning:
\begin{enumerate}
\item Initialise the surviving set $\mathcal{C}_0 = \{1,\ldots,K\}$.
\item For each tier $\ell = 0,1,2,3$ in priority order:
  \begin{enumerate}
  \item Compute $S_\ell^* = \min_{k \in \mathcal{C}_\ell} S_\ell^{(k)}$.
  \item Retain only near-optimal candidates:
    $\mathcal{C}_{\ell+1} = \{k \in \mathcal{C}_\ell : S_\ell^{(k)} \le S_\ell^* + \varepsilon_\ell\}$.
  \end{enumerate}
\item Among the survivors $\mathcal{C}_4$, select
  $k^* = \arg\min_{k \in \mathcal{C}_4}\sum_{\ell=0}^{3} S_\ell^{(k)}$.
\item If an exact tie remains after this residual aggregate comparison, return the lowest candidate index in the tied set.
\end{enumerate}
Return $\hat{\tau} = \tau^{(k^*)}$.
\end{definition}

\noindent
The definition is deliberately \emph{procedural} rather than order-theoretic. The tolerance relation $|S_\ell^{(i)} - S_\ell^{(j)}| \le \varepsilon_\ell$ is \emph{not} transitive, and consequently the pruning sets $\mathcal{C}_\ell$ depend on the global minimum within the current surviving pool, not on pairwise comparisons.  This means the procedure does not induce a total order on $\mathbb{R}_{\ge 0}^4$ in the classical sense; it defines a deterministic selection \emph{algorithm} rather than a classical preference relation.  As shown in Proposition~\ref{thm:order_invariance}, the output is nonetheless well-defined and invariant to candidate enumeration order.

For conciseness, we sometimes write $\mathbf{S}^{(i)} \prec_{\text{lex},\epsilon} \mathbf{S}^{(j)}$ to mean that candidate~$i$ would survive and candidate~$j$ would not whenever both are present in the same candidate pool---but we emphasise that this notation is a shorthand for the procedural definition above, not a claim that $\prec_{\text{lex},\epsilon}$ is a partial order.

Lexicographic selection requires at most $O(KL)$ tier comparisons (with $L{=}4$) and permits early termination: as soon as a candidate is pruned at a higher-priority tier, lower tiers cannot recover it.  The residual aggregate tie-break is consulted only after tolerance pruning has exhausted the strict tier ordering, so it resolves near-ties within the surviving set rather than redefining the priority hierarchy.

%-------------------------------------------------------------------------------
\subsection{Properties of Lexicographic Selection}
\label{subsec:properties}
%-------------------------------------------------------------------------------

Before proving formal guarantees, it is useful to state the qualitative properties that follow directly from the ordering definition, because they clarify why lexicographic selection matches traffic-rule semantics better than scalarization.

\begin{itemize}
\item \textbf{Priority preservation.}  If two candidates differ beyond tolerance at tier~$\ell$, then their ordering is decided at tier~$\ell$, independent of all lower tiers.

\item \textbf{No weight-induced trade-offs.}  Unlike weighted sums, the method does not require tuning cross-tier weights that implicitly allow comfort improvements to override safety or legality.

\item \textbf{Determinism with controlled tie-handling.}  The tolerance comparator reduces unstable ordering changes under small numerical perturbations; residual ties are resolved by aggregate residual score over the surviving set, with candidate index used only as a final exact-tie fallback.

\item \textbf{Auditability.}  The selected decision can be attributed to the first tier~$\ell^\star$ that differentiates candidates, enabling concise explanations (e.g., ``rejected due to Legal-tier red-light violation'').
\end{itemize}

\noindent
These properties support the central goal of this paper: selection should respect non-negotiable traffic priorities while remaining compatible with modern multi-modal generators.  The following subsection strengthens these qualitative claims with three formal theorems.

%-------------------------------------------------------------------------------
\subsection{Formal Guarantees}
\label{subsec:formal_guarantees}
%-------------------------------------------------------------------------------

We now establish three formal properties of the filter-based lexicographic selector (Algorithm~\ref{alg:lexicographic} in Section~\ref{Sec:RECTOR}). \emph{Order invariance}: the selection is a function of the candidate set, independent of enumeration order. \emph{Infeasibility fallback}: when no fully compliant candidate exists, the system performs graceful degradation by selecting the least-violating option rather than reverting to confidence-only selection. \emph{scalarization correctness}: the GPU-friendly scoring mechanism preserves the intended lexicographic ordering, ensuring lower-priority gains never override higher-priority requirements. Together, these results formalize the structural distinction introduced in Section~\ref{Sec:Introduction}: \textsc{RECTOR}'s lexicographic ordering \emph{provably cannot} trade higher-priority violations for lower-priority gains, whereas a weighted-sum baseline merely \emph{does not} under its current parameterization.

\begin{proposition}[Order Invariance]
\label{thm:order_invariance}
Algorithm~\ref{alg:lexicographic} produces the same output regardless of the ordering of candidates in the input set $\mathcal{T}$, for any fixed tolerance vector $(\varepsilon_0,\ldots,\varepsilon_3)$ with $\varepsilon_\ell \geq 0$.
\end{proposition}
\begin{proof}
At each tier~$\ell$, the algorithm computes $S^* = \min_{k \in \mathcal{C}} S_\ell^{(k)}$ and retains $\mathcal{C}' = \{k \in \mathcal{C} : S_\ell^{(k)} \leq S^* + \varepsilon_\ell\}$. Both the minimum operation and the threshold filter depend only on the \emph{set of tier scores}, not on the order in which candidates are enumerated. The residual aggregate tie-break $\arg\min_{k \in \mathcal{C}}\sum_{\ell} S_\ell^{(k)}$ is also order-independent, and the final exact-tie fallback uses a fixed deterministic rule (lowest candidate index). Therefore the output is invariant to input permutation. \qed
\end{proof}

The first proposition ensures that the selection is a well-defined function of the candidate \emph{set}, not an artifact of enumeration order.

\begin{remark}[Scope of order invariance]
\label{rem:order_invariance_scope}
Order invariance guarantees that Algorithm~\ref{alg:lexicographic} produces the same selected index $k^*$ regardless of the enumeration order of the $K$ candidates.  It does \emph{not} assert that $k^*$ is the unique global optimum under a transitive total order.  Because the tolerance relation $S_\ell^{(i)} \approx_\varepsilon S_\ell^{(j)}$ (Section~\ref{subsec:lexicographic_selection}) is \emph{not} transitive, the output depends on which candidates survive each tier's pruning set $\mathcal{C}_\ell$; this set is determined by the minimum tier score in the current candidate pool, not by pairwise comparisons.  In practice, with $\varepsilon_\ell \leq 10^{-3}$ and $K \leq 12$ (the main paper uses $K{=}6$), any non-transitive chain can involve only a small finite subset of candidates, and surviving candidates are typically numerically close enough that the residual aggregate tie-break resolves them deterministically.
\end{remark}

The second proposition addresses the practically important case in which the candidate set contains no fully compliant option.

\begin{proposition}[Infeasibility Fallback]
\label{thm:infeasibility}
When all $K$ candidates violate Safety-tier rules (i.e., $S_0^{(k)} > 0$ for all $k \in [K]$), Algorithm~\ref{alg:lexicographic} selects the candidate with minimum Safety-tier score within tolerance $\varepsilon_0$, breaking ties via lower tiers and then confidence. Formally, the algorithm returns the \emph{least-violating} candidate under the full lexicographic ordering.
\end{proposition}
\begin{proof}
The filter at tier~0 retains all candidates within $\varepsilon_0$ of $\min_k S_0^{(k)}$. Subsequent tiers further discriminate among these survivors. The procedure terminates with the candidate that minimizes the tier score vector lexicographically, which is, by construction, the least-violating trajectory. \qed
\end{proof}

\noindent\textbf{Remark (Infeasibility signalling).} When the selected trajectory has $S_0 > 0$ (Safety-tier violation), the system detects infeasibility.  The implementation now emits an explicit \texttt{infeasible} flag when the selected candidate has non-zero Safety-tier score, enabling downstream fail-safe logic such as emergency braking or conservative fallback.

The first two propositions address the selector's logical correctness.  The third shows that the lexicographic ordering can be implemented efficiently on GPU hardware via a scalar surrogate, without sacrificing correctness.

\begin{proposition}[scalarization Correctness]
\label{thm:scalarization}
The GPU-friendly scalarization $\mathrm{Score}^{(k)} = \sum_{\ell=0}^{3} B^{4-\ell} S_\ell^{(k)}$ with base $B$ is order-preserving with respect to exact lexicographic ordering (i.e., $\varepsilon_\ell = 0$ for all $\ell$) if and only if $\max_{k,\ell} S_\ell^{(k)} < B/(L-1)$.
\end{proposition}
\begin{proof}
For exact lexicographic ordering, candidate $i$ is preferred to $j$ iff there exists $\ell^*$ such that $S_\ell^{(i)} = S_\ell^{(j)}$ for $\ell < \ell^*$ and $S_{\ell^*}^{(i)} < S_{\ell^*}^{(j)}$. Under the scalarization, the difference at tier $\ell^*$ contributes at least $B^{4-\ell^*}$ to the score gap, while all lower tiers together contribute at most $(L - \ell^* - 1) \cdot \max_\ell S_\ell \cdot B^{3-\ell^*}$. This is dominated whenever $(L-1) \cdot \max_\ell S_\ell < B$, i.e.\ $\max_\ell S_\ell < B/(L-1)$.

In the RECTOR implementation, each tier score $S_\ell^{(k)}$ is a \emph{normalized weighted average} of per-rule severities: $S_\ell^{(k)} = \sum_{r:\ell(r)=\ell} \tilde{w}_r \bar{V}_r^{(k)}$ where the intra-tier weights satisfy $\tilde{w}_r \ge 0$, $\sum_r \tilde{w}_r = 1$, and each normalized violation $\bar{V}_r^{(k)} \in [0,1]$ (see Section~\ref{subsec:proxy_details}). Consequently, $\max_{k,\ell} S_\ell^{(k)} \leq 1$ regardless of the number of rules in any tier---in particular, regardless of the 14 rules in the Comfort tier. With $B = 1000$ and $L = 4$ the sufficient condition becomes $\max_{k,\ell} S_\ell^{(k)} \leq 1 < 1000/3 \approx 333$, which holds; any $B > L - 1 = 3$ would suffice. \qed
\end{proof}

\begin{corollary}
The bound $\max_\ell S_\ell < B/(L-1)$ is independent of the number of rules per tier as long as tier scores are normalized to $[0,1]$. A tier with 14 rules does not weaken the guarantee; only the normalization contract must be maintained by the proxy aggregator.
\end{corollary}

\begin{remark}[Normalization contract and choice of $B$]
\label{rem:B_choice}
Because tier scores are normalized to $[0,1]$ by the intra-tier weight constraint $\sum_{r:\ell(r)=\ell}\tilde{w}_r = 1$ (Section~\ref{subsec:tier_aggregation}), the sufficient condition $\max_{k,\ell} S_\ell^{(k)} \leq 1 < B/(L{-}1)$ is automatically satisfied for \emph{any} $B > L{-}1 = 3$.  In other words, \textbf{Proposition~\ref{thm:scalarization} is a straightforward consequence of the normalization contract}: given that tier scores lie in $[0,1]$ by construction, the scalarization is order-preserving for any base $B > 3$, and no empirical verification or careful parameter choice is required.  The reference implementation uses $B{=}1000$ for a single practical reason: it provides a $1000\times$ separation between adjacent tiers, giving numerical headroom for floating-point accumulation over $K \leq 6$ candidates in batched GPU execution.  The maximum composite score ($B^4 = 10^{12}$) remains well within IEEE\,754 single-precision range (${\approx}\,3.4\times10^{38}$).  However, one subtlety deserves mention: with a Comfort-tier multiplier of $10^0{=}1$ and Safety-tier contributions up to $10^{12}$, the representable Comfort-tier resolution in the composite score is approximately $10^{-9}$ relative to the Safety contribution, which may cause correct Comfort-tier tie-breaking to be lost to floating-point cancellation when Safety-tier scores differ.  \textbf{The scalarization is therefore used only as a GPU pre-filter}; final selection always uses the procedural lexicographic algorithm (Algorithm~2), which compares tier scores sequentially and is immune to floating-point cancellation.  In practice, this is further mitigated by the tolerance-based pruning (Definition~\ref{def:lex_selection}), which resolves near-ties before the scalarization is consulted.
\end{remark}

\noindent
Taken together, these three results establish that lexicographic selection is well-defined (Proposition~\ref{thm:order_invariance}), gracefully degrading (Proposition~\ref{thm:infeasibility}), and efficiently implementable (Proposition~\ref{thm:scalarization}).  The ``cannot versus does not'' distinction follows directly: no re-parameterization of the lexicographic selector can cause it to trade a Safety-tier violation for a Comfort-tier gain, because the tier filtering is applied \emph{before} lower tiers are consulted.  A weighted-sum scalarization, by contrast, can always be reparameterized to produce such a trade-off, even if its current weights happen to avoid it.

%-------------------------------------------------------------------------------
\subsection{Comparison with Alternative Selection Strategies}
\label{subsec:strategy_comparison}
%-------------------------------------------------------------------------------

To make the formal properties above concrete, it is useful to compare lexicographic selection against two common alternatives that serve as baselines in our experiments.

\textbf{Confidence-only selection} chooses the most likely candidate:
\begin{equation}
\tau_{\text{conf}}=\arg\max_{\tau \in \mathcal{T}} p(\tau),
\end{equation}
which is computationally simple but does not encode acceptability.

\textbf{Weighted-sum selection} aggregates rule penalties into a single scalar:
\begin{equation}
\tau_{\text{ws}}=\arg\min_{\tau \in \mathcal{T}} \sum_{r \in \mathcal{R}} w_r\,\hat{a}_r(s)\,\bar{V}_r(\tau;s).
\end{equation}
All baselines use the same predicted applicability~$\hat{a}_r(s)$ as RECTOR to ensure a fair comparison.  While the weighted-sum approach is flexible, it introduces implicit trade-offs via~$w_r$, and small changes in weights can reorder solutions in undesirable ways, especially across tiers with different semantics.  Table~\ref{tab:selection_strategy_qualitative} summarizes these strategies and highlights the structural role of tiered ordering in RECTOR.

\begin{table*}
\centering
\caption{Selection Strategy Comparison (Qualitative)}
\label{tab:selection_strategy_qualitative}
\begin{tabular}{lccc}
\toprule
\textbf{Strategy} & \textbf{Selection Criterion} & \textbf{Key Limitation} & \textbf{Interpretability} \\
\midrule
Confidence Only & $\arg\max p^{(k)}$ & No safety or legality awareness & Low \\
Weighted Sum & $\arg\min \sum_r w_r a_r \bar{V}_r$ & Weight tuning induces trade-offs & Medium \\
Lexicographic (RECTOR) & $\arg\min \mathbf{S}(\tau;s)$ under $\prec_{\text{lex},\epsilon}$ & Requires tier structure and tolerances & High \\
\bottomrule
\end{tabular}
\end{table*}

%-------------------------------------------------------------------------------
\subsection{Full-Catalog Evaluation Context}
\label{subsec:fullcatalog_context}
%-------------------------------------------------------------------------------

Before stating the formal problem, one further methodological distinction must be made because it affects how results are interpreted throughout the paper.  In practice, not all rules have equally convenient smooth proxies, and some require semantic signals or discrete conditions.  We therefore distinguish between two evaluation modes:
\begin{enumerate}
\item \textbf{Proxy-based selection} (Protocol~A), where $\hat{a}_r(s)$ and smooth evaluators are used to rank candidates robustly across 24/28~rules, and \item \textbf{Full-catalog auditing} (Protocol~B), where the complete Waymo-style evaluator is applied to the selected trajectory using oracle applicability~$a_r(s)$ across all 28~rules.
\end{enumerate}

A key point is that human driving in the dataset is not violation-free under either protocol; therefore, the goal is not ``zero violations,'' but a consistent reduction in high-priority (Safety/Legal) violations while preserving competitive accuracy.  The proxy-based evaluation (24/28~rules) used for selection provides the appropriate like-for-like comparison across strategies.  When the full 28-rule evaluator is applied, absolute violation rates increase for all methods---including ground-truth human trajectories---due to the stricter evaluator and the inclusion of rules without proxies.  This methodological difference means that full-catalog absolute rates are not directly interpretable for safety comparisons between methods; see Section~\ref{Sec:Known_Gaps_and_Next_Steps_Toward_Deployment} for further discussion.

%-------------------------------------------------------------------------------
\subsection{Formal Problem Statement}
\label{subsec:problem_statement}
%-------------------------------------------------------------------------------

With all definitions, properties, and evaluation context established, we can now state the rule-aware selection problem concisely as the objective that RECTOR optimizes.

\begin{problem}[Rule-Aware Trajectory Selection]
\label{prob:selection}
The rule-aware selection problem is specified by a five-component tuple
$\langle f_\theta,\; \mathcal{R},\; \hat{a},\; \bar{V},\; \textsc{Select}\rangle$:
\begin{enumerate}[label=(\roman*),leftmargin=2em]
\item \textbf{Generator} $f_\theta$: maps scenario $s\sim\mathcal{D}$ to candidates and confidences $(\mathcal{T},\mathbf{p})$;
\item \textbf{Rule catalog} $\mathcal{R}$ with tier mapping
  $\ell\colon\mathcal{R}\to\{0,1,2,3\}$;
\item \textbf{Applicability predictor} $\hat{a}_r(s)\in\{0,1\}$: a learned neural head that predicts which rules are active in scenario~$s$
  (approximation source~1; see Section~\ref{Sec:RECTOR}); smooth proxies that approximate full rule-violation semantics
  (approximation source~2; see Section~\ref{Sec:Rule_Evaluation_And_Data_Pipeline});
\item \textbf{Selector} $\textsc{Select}$: tolerance-based lexicographic selection (Definition~\ref{def:lex_selection}) over tier scores $\mathbf{S}(\tau;s)$ aggregated from~(ii)--(iv).
\end{enumerate}
The objective is to select $\hat{\tau}\in\mathcal{T}$ that minimizes the tier score vector under the procedural lexicographic ordering:
\begin{equation}
\hat{\tau}=\textsc{Select}\!\big(\mathcal{T},\;\mathbf{S}(\cdot\,;s),\;\mathbf{p},\;\boldsymbol{\varepsilon}\big).
\end{equation}
\end{problem}

\noindent
The tuple formulation makes explicit that the deployed system involves two sources of approximation---applicability prediction~(iii) and proxy evaluation~(iv)---that can each introduce gaps relative to the idealised selection objective.  Correctness of the selector~(v) is established by Propositions~\ref{thm:order_invariance}--\ref{thm:scalarization} under the assumption that tier scores are exact; the empirical impact of the approximation gaps is quantified via the dual-protocol evaluation (Section~\ref{subsec:fullcatalog_context}) and the proxy-alignment analysis in Section~\ref{Sec:Rule_Evaluation_And_Data_Pipeline}.

This formulation makes the separation between \emph{generation} and \emph{selection} explicit and enables the central claim of this work: when multi-modal diversity exists, rule-aware selection can convert that diversity into measurable compliance gains.  In the present benchmark, the lexicographic selector contributes a formal priority guarantee and auditability rather than an empirical compliance separation from weighted-sum.  Section~\ref{Sec:RECTOR} describes the architecture that implements this objective, and Section~\ref{Sec:Accuracy_SectionRule-Compliance_Behavior_and_Efficiency}
presents the empirical evidence.